% Source: physical pp. 123--132 (printed pp. 100--109). Coverage: complete Section 17.5 through the paragraph before Problems; Eqs. (17.5.1)--(17.5.44), all unnumbered displays, Figure 17.1, and two footnotes. Uncertainties: none after full-resolution rendered-page review.
\subsection{A One--Loop Calculation in Background Field Gauge}

As an exercise, we are now going to calculate the one-loop renormalization
factor for the gauge coupling constant in a general non-Abelian gauge
theory. As we will see in the next chapter, this provides an essential input
in so-called `renormalization group' calculations of physical processes at
high energy; the results we obtain here will be used there to demonstrate
the asymptotic freedom of non-Abelian gauge theories.

The method to be employed here is somewhat novel. Usually, one
considers the effective action in a spacetime-dependent background gauge
field, and calculates the terms quadratic in this field, extracting a factor
$(q_\mu A_{\alpha\nu}-q_\nu A_{\alpha\mu})^2$ (where $q$ is the gauge field
four-momentum), and only then isolating the logarithmic divergence by setting
$q=0$ in the coefficient of this factor. Instead we shall follow the much
simpler course of taking the gauge field to be spacetime-independent from the
beginning. In this case the terms in the effective action that are quadratic
or cubic in the gauge field of course vanish, but there is a non-vanishing
quartic term that is ultraviolet divergent, and which can be used to calculate
the coupling constant renormalization factor $(1+L_A)^{-1/2}$. In this way our
one-loop calculation becomes a matter of simple matrix algebra. Note that
this procedure can only work in background field gauge; otherwise there
would be independent logarithmic divergences in the parts of the effective
action that are quartic and quadratic in the background gauge field.

With this motivation, we turn to the calculation of the one-loop effective
action in a background field for which $A_{\alpha\mu}$ is constant and
$\psi=\omega=\omega^*=0$. For such a background field, the full modified
Lagrangian is\footnote{See Eqs.~\eqref{eq:17.4.12}, \eqref{eq:17.4.4}, and
\eqref{eq:17.4.23}. We are here specializing to the case of matter fields
forming a multiplet of spin $1/2$ fermions. The squares in $\mathcal L$ and
$\mathcal L_f$ include obvious index contractions.}
\begin{equation}
\mathcal L_{\mathrm{MOD}}=\mathcal L+\mathcal L_f+\mathcal L_{GH},
\label{eq:17.5.1}
\end{equation}
\begin{align}
\mathcal L={}&-\frac14\left(F_{\alpha\mu\nu}+\bar D_\mu A'_{\alpha\nu}
-\bar D_\nu A'_{\alpha\mu}+C_{\alpha\beta\gamma}A_{\beta\mu}'A_{\gamma\nu}'\right)^2
\notag\\
&-\bar\psi'\bigl(\sl{\bar D}-it_\alpha\sl{A}'_\alpha+m\bigr)\psi',
\label{eq:17.5.2}
\end{align}
\begin{equation}
\mathcal L_f=-\frac{1}{2\xi}f_\alpha f_\alpha
=-\frac{1}{2\xi}\bigl(\bar D_\mu A_\alpha^{\prime\mu}\bigr)^2,
\label{eq:17.5.3}
\end{equation}
\begin{equation}
\mathcal L_{GH}=-\bigl(\bar D_\mu\omega_\alpha^{\prime *}\bigr)
\bigl(\bar D^\mu\omega'_\alpha
-C_{\alpha\beta\gamma}\omega'_\beta A_\gamma^{\prime\mu}\bigr).
\label{eq:17.5.4}
\end{equation}
One-loop graphs for vacuum--vacuum amplitudes are calculated from the
part of the action that is quadratic in the quantum fields
$A',\psi',\omega',\omega^{\prime *}$ over which one integrates. Keeping only
such quadratic terms, we have
\begin{align}
\mathcal L_{\mathrm{QUAD}}={}&-\frac14
\bigl(\bar D_\mu A'_{\alpha\nu}-\bar D_\nu A'_{\alpha\mu}\bigr)^2
-\frac12F_\alpha^{\mu\nu}C_{\alpha\beta\gamma}
A'_{\beta\mu}A'_{\gamma\nu}
\notag\\
&-\bar\psi'(\sl{\bar D}+m)\psi'
-\frac{1}{2\xi}\bigl(\bar D_\mu A_\alpha^{\prime\mu}\bigr)^2
-\bigl(\bar D_\mu\omega_\alpha^{\prime *}\bigr)
 \bigl(\bar D^\mu\omega'_\alpha\bigr).
\label{eq:17.5.5}
\end{align}
The corresponding action may be put in the general quadratic form:
\begin{align}
I_{\mathrm{QUAD}}\equiv{}&\int d^4x\,\mathcal L_{\mathrm{QUAD}}
\notag\\
={}&-\frac12\int d^4x\,d^4y\,
A_\alpha^{\prime\mu}(x)A_\beta^{\prime\nu}(y)
\mathcal D^A_{x\alpha\mu,y\beta\nu}
-\int d^4x\,d^4y\,\bar\psi'_k(x)\psi'_\ell(y)
\mathcal D^\psi_{xk,y\ell}
\notag\\
&-\int d^4x\,d^4y\,\omega_\alpha^{\prime *}(x)\omega'_\beta(y)
\mathcal D^\omega_{x\alpha,y\beta},
\label{eq:17.5.6}
\end{align}
with
\begin{align}
\mathcal D^A_{x\alpha\mu,y\beta\nu}\equiv{}&
\eta_{\mu\nu}
\left(-\delta_{\gamma\alpha}\frac{\partial}{\partial x_\lambda}
+C_{\gamma\delta\alpha}A_\delta^\lambda(x)\right)
\left(-\delta_{\gamma\beta}\frac{\partial}{\partial y^\lambda}
+C_{\gamma\epsilon\beta}A_{\epsilon\lambda}(y)\right)\delta^4(x-y)
\notag\\
&-\left(-\delta_{\gamma\alpha}\frac{\partial}{\partial x^\nu}
+C_{\gamma\delta\alpha}A_{\delta\nu}(x)\right)
\left(-\delta_{\gamma\beta}\frac{\partial}{\partial y^\mu}
+C_{\gamma\epsilon\beta}A_\epsilon^\mu(y)\right)\delta^4(x-y)
\notag\\
&+F_{\gamma\mu\nu}(x)C_{\gamma\alpha\beta}\delta^4(x-y)
\notag\\
&+\frac1\xi
\left(-\delta_{\gamma\alpha}\frac{\partial}{\partial x^\mu}
+C_{\gamma\delta\alpha}A_{\delta\mu}(x)\right)
\left(-\delta_{\gamma\beta}\frac{\partial}{\partial y^\nu}
+C_{\gamma\epsilon\beta}A_{\epsilon\nu}(y)\right)\delta^4(x-y),
\label{eq:17.5.7}
\end{align}
\begin{equation}
\mathcal D^\psi_{xk,y\ell}\equiv
\left(-\gamma^\mu\frac{\partial}{\partial y^\mu}
-it_\alpha\sl{A}_\alpha(y)+m\right)_{k\ell}\delta^4(x-y),
\label{eq:17.5.8}
\end{equation}
\begin{align}
\mathcal D^\omega_{x\alpha,y\beta}\equiv{}&
\left(-\delta_{\gamma\alpha}\frac{\partial}{\partial x^\lambda}
+C_{\gamma\delta\alpha}A_{\delta\lambda}(x)\right)
\left(-\delta_{\gamma\beta}\frac{\partial}{\partial y_\lambda}
+C_{\gamma\epsilon\beta}A_\epsilon^\lambda(y)\right)\delta^4(x-y).
\label{eq:17.5.9}
\end{align}
(The minus signs in front of $\partial/\partial x$ and $\partial/\partial y$
drop out when we integrate by parts.)

The one-loop contribution to the effective action is given (as in Section
16.2) by
\begin{align}
\exp\bigl(i\Gamma^{(1\ \mathrm{loop})}[A]\bigr)\propto{}&
\int_{1\mathrm{PI}}
\left(\prod dA'\right)\left(\prod d\psi'\right)
\left(\prod d\bar\psi'\right)\left(\prod d\omega'\right)
\left(\prod d\omega^{\prime *}\right)
\notag\\[-2pt]
&\qquad\cdot
\exp\bigl(iI_{\mathrm{QUAD}}[A',\psi',\bar\psi',\omega',
\omega^{\prime *};A]\bigr)
\notag\\
\propto{}&
\bigl(\operatorname{Det}\mathcal D^A\bigr)^{-1/2}
\bigl(\operatorname{Det}\mathcal D^\psi\bigr)^{+1}
\bigl(\operatorname{Det}\mathcal D^\omega\bigr)^{+1}.
\label{eq:17.5.10}
\end{align}
(The exponents $-1/2$ and $+1$ appear because $A'$ is a real boson field,
while $\psi'$, $\bar\psi'$, $\omega'$, and $\omega^{\prime *}$ are distinct
fermionic fields.) The calculation of such determinants is generally not
an easy task. However, it becomes much simpler in the case of constant
external fields, where the $\mathcal D$s can be diagonalized by passing to
momentum space.

Let us therefore now consider the case of constant background field
$A_{\alpha\mu}$. For non-Abelian gauge theories such a constant field
\emph{cannot} be removed by a gauge transformation, as shown by the non-zero
values of various gauge-covariant fields
\begin{equation}
F_{\alpha\mu\nu}=C_{\alpha\beta\gamma}A_{\beta\mu}A_{\gamma\nu},
\label{eq:17.5.11}
\end{equation}
\begin{equation}
D_\lambda F_{\delta\mu\nu}
=C_{\delta\epsilon\alpha}C_{\alpha\beta\gamma}
A_{\epsilon\lambda}A_{\beta\mu}A_{\gamma\nu},
\label{eq:17.5.12}
\end{equation}
and so on. Lorentz and gauge invariance tell us how to express the part
of $\Gamma[A]$ of a given dimensionality as the integral of a finite number
of local functions of $F_{\alpha\mu\nu}$, $D_\lambda F_{\alpha\mu\nu}$, etc.;
the coefficients of the terms in this expression can be inferred by comparing
the contribution that these terms make to $\Gamma[A]$ for constant background
field $A_{\alpha\mu}$ with the results of a perturbative expansion.

We transform each of the `matrices' $\mathcal D^A$, $\mathcal D^\psi$, and
$\mathcal D^\omega$ to a momentum basis by the usual normalized Fourier transform
\begin{equation}
\mathcal D_{q\cdots,p\cdots}
=\int\frac{d^4x}{(2\pi)^2}e^{-iq\cdot x}
 \int\frac{d^4y}{(2\pi)^2}e^{ip\cdot y}\mathcal D_{x\cdots,y\cdots}.
\label{eq:17.5.13}
\end{equation}
With $A$ constant, this gives
\begin{equation}
\mathcal D_{q\cdots,p\cdots}=\delta^4(p-q)\mathcal M_{\cdots,\cdots}(q),
\label{eq:17.5.14}
\end{equation}
where $\cdots$ denotes discrete indices, and the $\mathcal M$ are finite
$q$-dependent matrices
\begin{align}
\mathcal M^A_{\alpha\mu,\beta\nu}(q)={}&
\eta_{\mu\nu}\bigl(-iq^\lambda\delta_{\gamma\alpha}
+A_\delta^\lambda C_{\gamma\delta\alpha}\bigr)
\bigl(iq_\lambda\delta_{\gamma\beta}
+A_{\epsilon\lambda}C_{\gamma\epsilon\beta}\bigr)
\notag\\
&-\bigl(-iq_\nu\delta_{\gamma\alpha}
+A_{\delta\nu}C_{\gamma\delta\alpha}\bigr)
\bigl(iq_\mu\delta_{\gamma\beta}
+A_\epsilon^\mu C_{\gamma\epsilon\beta}\bigr)
+F_{\gamma\mu\nu}C_{\gamma\alpha\beta}
\notag\\
&+\frac1\xi\bigl(-iq_\mu\delta_{\gamma\alpha}
+A_{\delta\mu}C_{\gamma\delta\alpha}\bigr)
\bigl(iq_\nu\delta_{\gamma\beta}
+A_{\epsilon\nu}C_{\gamma\epsilon\beta}\bigr)
+\epsilon\text{ terms},
\label{eq:17.5.15}
\end{align}
\begin{equation}
\mathcal M^\psi_{k\ell}(q)
=\bigl(i\sl q-it_\alpha\sl A_\alpha+m\bigr)_{k\ell}
+\epsilon\text{ terms},
\label{eq:17.5.16}
\end{equation}
\begin{align}
\mathcal M^\omega_{\alpha\beta}(q)={}&
\bigl(-iq_\lambda\delta_{\gamma\alpha}
+A_{\delta\lambda}C_{\gamma\delta\alpha}\bigr)
\bigl(iq^\lambda\delta_{\gamma\beta}
+A_\epsilon^\lambda C_{\gamma\epsilon\beta}\bigr)
+\epsilon\text{ terms},
\label{eq:17.5.17}
\end{align}
with $F_{\gamma\mu\nu}$ given by Eq.~\eqref{eq:17.5.11}. From
Eq.~\eqref{eq:17.5.10} we have then
\begin{align}
i\Gamma^{(1\ \mathrm{loop})}[A]
={}&-\frac12\ln\operatorname{Det}\mathcal D^A
+\ln\operatorname{Det}\mathcal D^\psi
+\ln\operatorname{Det}\mathcal D^\omega
\notag\\
={}&-\frac12\operatorname{Tr}\ln\mathcal D^A
+\operatorname{Tr}\ln\mathcal D^\psi
+\operatorname{Tr}\ln\mathcal D^\omega
\notag\\
={}&\delta^4(p-p)\int d^4q\left[-\frac12\operatorname{tr}\ln\mathcal M^A(q)
+\operatorname{tr}\ln\mathcal M^\psi(q)
+\operatorname{tr}\ln\mathcal M^\omega(q)\right].
\label{eq:17.5.18}
\end{align}
We denote traces by `tr' instead of `Tr' in the last line of
Eq.~\eqref{eq:17.5.18} (and from now on in this section) to indicate that
these are the usual traces of finite matrices rather than of integral operators.

Since we are aiming here at a calculation of the infinite factor $L_A$
multiplying $FF$ terms in the effective action, let's isolate the term in
\eqref{eq:17.5.18} that is of fourth order in the background field $A$. For
this purpose, it is convenient to divide each $\mathcal M$ into terms
$\mathcal M_n$ containing $n=0,1$, or 2 factors of $A$:
\begin{equation}
\mathcal M=\mathcal M_0+\mathcal M_1+\mathcal M_2.
\label{eq:17.5.19}
\end{equation}
It is then elementary algebra to show that the term in
\eqref{eq:17.5.18} of fourth order in $A_{\alpha\mu}$ is

% Source: physical p. 127 (printed p. 104). Coverage: complete Figure 17.1 and caption. Uncertainty: none; geometry reconstructed from the rendered source.
\begin{figure}[htbp]
  \centering
  \begin{tikzpicture}[
    x=1cm,y=1cm,
    internal/.style={draw,line width=0.65pt},
    background/.style={draw,dashed,line width=0.65pt,dash pattern=on 3pt off 2.4pt}
  ]
    % Two-vertex loop. Each vertex carries two background-field factors.
    \draw[internal] (0,0) circle[radius=0.62];
    \draw[background] (-0.62,0) -- (-1.45,0.48);
    \draw[background] (-0.62,0) -- (-1.45,-0.48);
    \draw[background] (0.62,0) -- (1.45,0.48);
    \draw[background] (0.62,0) -- (1.45,-0.48);

    % Triangular loop. Two background factors meet the left vertex, one each
    % meets the upper-right and lower-right vertices.
    \coordinate (tL) at (3.35,0);
    \coordinate (tU) at (4.55,0.72);
    \coordinate (tD) at (4.55,-0.72);
    \draw[internal] (tL) -- (tU) -- (tD) -- cycle;
    \draw[background] (tL) -- (2.48,0.48);
    \draw[background] (tL) -- (2.48,-0.48);
    \draw[background] (tU) -- (5.38,1.22);
    \draw[background] (tD) -- (5.38,-1.22);

    % Four-vertex square loop. Each corner carries one background factor.
    \coordinate (sUL) at (7.25,0.62);
    \coordinate (sUR) at (8.35,0.62);
    \coordinate (sLR) at (8.35,-0.62);
    \coordinate (sLL) at (7.25,-0.62);
    \draw[internal] (sUL) -- (sUR) -- (sLR) -- (sLL) -- cycle;
    \draw[background] (sUL) -- (6.48,1.39);
    \draw[background] (sUR) -- (9.12,1.39);
    \draw[background] (sLR) -- (9.12,-1.39);
    \draw[background] (sLL) -- (6.48,-1.39);
  \end{tikzpicture}
  \renewcommand{\thefigure}{17.1}
  \caption{One-loop Feynman diagrams for the term in the quantum effective action that is quartic in a constant background gauge field $A^\mu_\alpha$. Here solid lines represent internal gauge, ghost, or matter lines; dashed lines indicate factors of $A^\mu_\alpha$. These three diagrams correspond to the three terms in Eq.~\eqref{eq:17.5.20}.}
  \label{fig:17.1}
\end{figure}


\begin{equation}
\left[\operatorname{tr}\ln\mathcal M\right]_{A^4}
=\operatorname{tr}\left\{-\frac12\left[\mathcal M_0^{-1}\mathcal M_2\right]^2
+\left[\mathcal M_0^{-1}\mathcal M_1\right]^2\mathcal M_0^{-1}\mathcal M_2
-\frac14\left[\mathcal M_0^{-1}\mathcal M_1\right]^4\right\}.
\label{eq:17.5.20}
\end{equation}
(To see this, insert factors $\epsilon$ and $\epsilon^2$ multiplying
$\mathcal M_1$ and $\mathcal M_2$ in Eq.~\eqref{eq:17.5.19}, differentiate
$\operatorname{tr}\ln\mathcal M$ four times with respect to $\epsilon$, divide
by $4!$ and set $\epsilon=0$.) The $\mathcal M_0^{-1}$ factors here are just
the usual propagators; for $\xi=1$, these are
\begin{equation}
\left[\mathcal M_0^A(q)\right]^{-1}_{\alpha\mu,\beta\nu}
=\delta_{\alpha\beta}\eta_{\mu\nu}(q^2-i\epsilon)^{-1},
\label{eq:17.5.21}
\end{equation}
\begin{equation}
\left[\mathcal M_0^\psi(q)\right]^{-1}_{k\ell}
=\left[i\sl q+m\right]^{-1}_{k\ell},
\label{eq:17.5.22}
\end{equation}
\begin{equation}
\left[\mathcal M_0^\omega(q)\right]^{-1}_{\alpha,\beta}
=\delta_{\alpha\beta}(q^2-i\epsilon)^{-1}.
\label{eq:17.5.23}
\end{equation}
Indeed, the three terms in Eq.~\eqref{eq:17.5.20} just correspond to the
three Feynman diagrams shown in Figure~\ref{fig:17.1}; the present method
of calculation saves us from having to think about signs and combinatoric
factors.

For the $A$ loop, Eq.~\eqref{eq:17.5.15} gives for $\xi=1$:
\begin{align*}
\left[\mathcal M_1^A(q)\right]_{\alpha\mu,\beta\nu}
&=-2\eta_{\mu\nu}q^\lambda[\mathcal A_\lambda]_{\alpha\beta},\\
\left[\mathcal M_2^A(q)\right]_{\alpha\mu,\beta\nu}
&=\left[\eta_{\mu\nu}\mathcal A_\lambda\mathcal A^\lambda
-\mathcal A_\nu\mathcal A_\mu-\mathcal A_\mu\mathcal A_\nu\right]_{\alpha\beta}
+F_{\gamma\mu\nu}C_{\gamma\alpha\beta},
\end{align*}
where $\mathcal A_\lambda$ is the matrix
\[
[\mathcal A_\lambda]_{\alpha\beta}\equiv-iC_{\alpha\beta\gamma}A_{\gamma\lambda}
\]
for which
\begin{align*}
[\mathcal A_\lambda,\mathcal A_\rho]_{\alpha\beta}
&=-C_{\alpha\delta\gamma}C_{\delta\beta\epsilon}
(A_{\gamma\lambda}A_{\epsilon\rho}-A_{\gamma\rho}A_{\epsilon\lambda})\\
&=-\bigl(C_{\alpha\delta\gamma}C_{\delta\beta\epsilon}
+C_{\alpha\delta\epsilon}C_{\delta\beta\gamma}\bigr)
A_{\gamma\lambda}A_{\epsilon\rho}\\
&=+C_{\alpha\delta\beta}C_{\delta\epsilon\gamma}
A_{\gamma\lambda}A_{\epsilon\rho}
=C_{\alpha\delta\beta}F_{\delta\rho\lambda}.
\end{align*}
The integrals have the structure
\begin{align*}
\int d^4q\,q^\mu q^\nu f(q^2)
&=\frac14\eta^{\mu\nu}\int d^4q\,q^2f(q^2),\\
\int d^4q\,q^\mu q^\nu q^\rho q^\sigma f(q^2)
&=\frac1{24}\bigl[\eta^{\mu\nu}\eta^{\rho\sigma}
+\eta^{\mu\rho}\eta^{\nu\sigma}
+\eta^{\mu\sigma}\eta^{\nu\rho}\bigr]
\int d^4q\,(q^2)^2f(q^2).
\end{align*}
We then find, for $\xi=1$:
\begin{align*}
\int d^4q\,\operatorname{tr}\left\{
\left[\mathcal M_0^A(q)^{-1}\mathcal M_2^A(q)\right]^2\right\}
\notag\\[-2pt]
={}&4\mathcal I\,\operatorname{tr}
[\mathcal A^\lambda\mathcal A_\lambda\mathcal A^\eta\mathcal A_\eta]
+4\mathcal I\,C_{\gamma\alpha\beta}C_{\delta\alpha\beta}
F_{\gamma\mu\nu}F_\delta^{\mu\nu},\\
\int d^4q\,\operatorname{tr}\left\{
\left[\mathcal M_0^A(q)^{-1}\mathcal M_1^A(q)\right]^2
\mathcal M_0^A(q)^{-1}\mathcal M_2^A(q)\right\}
\notag\\[-2pt]
={}&4\mathcal I\,\operatorname{tr}
[\mathcal A^\lambda\mathcal A_\lambda\mathcal A^\eta\mathcal A_\eta],\\
\int d^4q\,\operatorname{tr}\left\{
\left[\mathcal M_0^A(q)^{-1}\mathcal M_1^A(q)\right]^4\right\}
\notag\\[-2pt]
={}&\frac83\mathcal I\,\operatorname{tr}
\left[\begin{aligned}
&2\mathcal A^\lambda\mathcal A_\lambda\mathcal A^\eta\mathcal A_\eta\\[-2pt]
&+\mathcal A^\lambda\mathcal A^\eta\mathcal A_\lambda\mathcal A_\eta
\end{aligned}\right],
\end{align*}
where $\mathcal I$ is the divergent integral
\begin{equation}
\mathcal I\equiv\int d^4q\,[q^2-i\epsilon]^{-2},
\label{eq:17.5.24}
\end{equation}
whose significance is discussed below. Putting this together in
Eq.~\eqref{eq:17.5.20}, we have
\begin{align*}
\int d^4q\left[\operatorname{tr}\ln\mathcal M^A(q)\right]_{A^4}
={}&\frac23\mathcal I\,\operatorname{tr}
[\mathcal A^\lambda\mathcal A_\lambda\mathcal A^\eta\mathcal A_\eta
-\mathcal A^\lambda\mathcal A^\eta\mathcal A_\lambda\mathcal A_\eta]
\notag\\[-2pt]
&-2\mathcal I\,C_{\gamma\alpha\beta}C_{\delta\alpha\beta}
F_{\gamma\mu\nu}F_\delta^{\mu\nu}.
\end{align*}
Both terms are actually of the same form, and when combined yield
\begin{equation}
\int d^4q\left[\operatorname{tr}\ln\mathcal M^A(q)\right]_{A^4}
=-\frac53\mathcal I\,C_{\gamma\alpha\beta}C_{\delta\alpha\beta}
F_{\gamma\mu\nu}F_\delta^{\mu\nu}.
\label{eq:17.5.25}
\end{equation}
Now jumping ahead to the ghost loop, we see from Eq.~\eqref{eq:17.5.17} that
\begin{equation}
[\mathcal M_1^\omega(q)]_{\alpha\beta}
=-2[\mathcal A^\lambda]_{\alpha\beta}q_\lambda,
\label{eq:17.5.26}
\end{equation}
\begin{equation}
[\mathcal M_2^\omega(q)]_{\alpha\beta}
=[\mathcal A^\lambda\mathcal A_\lambda]_{\alpha\beta}.
\label{eq:17.5.27}
\end{equation}
We then find
\begin{align*}
\int d^4q\,\operatorname{tr}\left\{
[\mathcal M_0^\omega(q)^{-1}\mathcal M_2^\omega(q)]^2\right\}
\notag\\[-2pt]
&=\mathcal I\,\operatorname{tr}
[\mathcal A^\lambda\mathcal A_\lambda\mathcal A^\eta\mathcal A_\eta],\\
\int d^4q\,\operatorname{tr}\left\{
[\mathcal M_0^\omega(q)^{-1}\mathcal M_1^\omega(q)]^2
\mathcal M_0^\omega(q)^{-1}\mathcal M_2^\omega(q)\right\}
\notag\\[-2pt]
&=\mathcal I\,\operatorname{tr}
[\mathcal A^\lambda\mathcal A_\lambda\mathcal A^\eta\mathcal A_\eta],\\
\int d^4q\,\operatorname{tr}\left\{
[\mathcal M_0^\omega(q)^{-1}\mathcal M_1^\omega(q)]^4\right\}
\notag\\[-2pt]
&=\frac23\mathcal I\,\operatorname{tr}
\left[\begin{aligned}
&2\mathcal A^\lambda\mathcal A_\lambda\mathcal A^\eta\mathcal A_\eta\\[-2pt]
&+\mathcal A^\lambda\mathcal A^\eta\mathcal A_\lambda\mathcal A_\eta
\end{aligned}\right].
\end{align*}
Thus for the ghost loop the integral of the quantity \eqref{eq:17.5.20} is
\begin{align}
\int d^4q\left[\operatorname{Tr}\ln\mathcal M^\omega(q)\right]_{A^4}
={}&\frac16\mathcal I\,\operatorname{Tr}
[\mathcal A^\lambda\mathcal A_\lambda\mathcal A^\eta\mathcal A_\eta
-\mathcal A^\lambda\mathcal A^\eta\mathcal A_\lambda\mathcal A_\eta]
\notag\\
={}&\frac1{12}\mathcal I\,C_{\gamma\alpha\beta}C_{\delta\alpha\beta}
F_{\gamma\mu\nu}F_\delta^{\mu\nu}.
\label{eq:17.5.28}
\end{align}
Finally, the vertices in the matter loop are
\[
[\mathcal M_1^\psi(q)]_{k\ell}=-i(t_\alpha\sl A_\alpha)_{k\ell},
\qquad
[\mathcal M_2^\psi(q)]_{k\ell}=0,
\]
so here there is only one term in Eq.~\eqref{eq:17.5.20}:
\[
\int d^4q\left[\operatorname{Tr}\ln\mathcal M^\psi(q)\right]_{A^4}
=-\frac14\int d^4q\,\operatorname{Tr}
\left\{\left[(i\sl q+m)^{-1}t_\alpha\sl A_\alpha\right]^4\right\}.
\]
We are interested in the ultraviolet-divergent part of this integral, so we
may drop the mass (which is negligible for large $q^\nu$) and write
\begin{align}
\int d^4q\left[\operatorname{Tr}\ln\mathcal M^\psi(q)\right]_{A^4}
={}&-\frac14\operatorname{Tr}\{t_\alpha t_\beta t_\gamma t_\delta\}
A_{\alpha\mu}A_{\beta\nu}A_{\gamma\rho}A_{\delta\sigma}
\notag\\[-2pt]
&\quad\cdot
\int d^4q\,
\frac{\operatorname{Tr}\{\sl q\gamma^\mu\sl q\gamma^\nu
\sl q\gamma^\rho\sl q\gamma^\sigma\}}{(q^2-i\epsilon)^8}
\notag\\
={}&-\frac{\mathcal I}{96}\operatorname{Tr}\{t_\alpha t_\beta t_\gamma t_\delta\}
A_{\alpha\mu}A_{\beta\nu}A_{\gamma\rho}A_{\delta\sigma}
\notag\\
&\quad\cdot\operatorname{Tr}\left\{
\begin{aligned}
&2\gamma_\lambda\gamma^\mu\gamma^\lambda\gamma^\nu
 \gamma_\eta\gamma^\rho\gamma^\eta\gamma^\sigma\\[-2pt]
&+\gamma_\lambda\gamma^\mu\gamma_\eta\gamma^\nu
 \gamma^\lambda\gamma^\rho\gamma^\eta\gamma^\sigma
\end{aligned}\right\}.
\label{eq:17.5.29}
\end{align}
where $\mathcal I$ is the same divergent integral as in Eq.~\eqref{eq:17.5.24}.
To calculate the traces of Dirac matrices, we use the anticommutation
relations of these matrices to write
\begin{align*}
\operatorname{Tr}\left\{\begin{aligned}
&2\gamma_\lambda\gamma^\mu\gamma^\lambda\gamma^\nu
 \gamma_\eta\gamma^\rho\gamma^\eta\gamma^\sigma\\[-2pt]
&+\gamma_\lambda\gamma^\mu\gamma_\eta\gamma^\nu
 \gamma^\lambda\gamma^\rho\gamma^\eta\gamma^\sigma
\end{aligned}\right\}
={}&8\operatorname{Tr}\{\gamma^\mu\gamma^\nu\gamma^\rho\gamma^\sigma\}
\notag\\[-2pt]
&\quad-4\operatorname{Tr}\{\gamma^\nu\gamma^\mu\gamma^\rho\gamma^\sigma\}
-4\operatorname{Tr}\{\gamma^\mu\gamma^\rho\gamma^\nu\gamma^\sigma\}\\
={}&-64\eta^{\mu\rho}\eta^{\nu\sigma}
+32\eta^{\mu\nu}\eta^{\rho\sigma}
+32\eta^{\mu\sigma}\eta^{\nu\rho}.
\end{align*}
Eq.~\eqref{eq:17.5.29} then gives
\begin{align}
\int d^4q\left[\operatorname{Tr}\ln\mathcal M^\psi(q)\right]_{A^4}
={}&\frac13\mathcal I\,\operatorname{Tr}
\{[t_\alpha,t_\beta][t_\gamma,t_\delta]\}
A_{\alpha\mu}A_{\beta\nu}A_\gamma^\mu A_\delta^\nu
\notag\\
={}&-\frac13\mathcal I\,F_{\gamma\mu\nu}F_\delta^{\mu\nu}
\operatorname{Tr}\{t_\gamma,t_\delta\}.
\label{eq:17.5.30}
\end{align}
Using Eqs.~\eqref{eq:17.5.25}, \eqref{eq:17.5.28}, and
\eqref{eq:17.5.30} in Eq.~\eqref{eq:17.5.18} gives at last
\begin{equation}
\Gamma_{A^4}^{(1\ \mathrm{loop})}
=\frac{-i\mathcal I}{(2\pi)^4}\int d^4x\,F_{\gamma\mu\nu}F_\delta^{\mu\nu}
\left[\left(\frac56+\frac1{12}\right)
C_{\gamma\alpha\beta}C_{\delta\alpha\beta}
-\frac13\operatorname{Tr}\{t_\gamma,t_\delta\}\right],
\label{eq:17.5.31}
\end{equation}
where we have expressed the momentum-space delta-function in
\eqref{eq:17.5.14} as
\begin{equation}
\delta^4(p-p)=(2\pi)^{-4}\int d^4x\,1.
\label{eq:17.5.32}
\end{equation}
It is important that the result turns out to depend on $A_{\alpha\mu}$ only
through the field strength \eqref{eq:17.5.11}, as required by background
gauge invariance.

Let us now (for the first time in this section) use the assumed simplicity
of the gauge group and irreducibility of the matter field multiplet. In this
case
\begin{equation}
C_{\gamma\alpha\beta}C_{\delta\alpha\beta}=g^2C_1\delta_{\gamma\delta},
\label{eq:17.5.33}
\end{equation}
\begin{equation}
\operatorname{Tr}\{t_\gamma,t_\delta\}=g^2C_2\delta_{\gamma\delta},
\label{eq:17.5.34}
\end{equation}
where $g$ is the common gauge coupling constant appearing as a factor in
$C_{\gamma\alpha\beta}$ and $t_\gamma$, and $C_1$ and $C_2$ are numerical
constants that characterize the gauge group and the representation of this
group provided by the matter multiplet. For instance, in the original
Yang--Mills theory the gauge group is $SU(2)$ (or equivalently $SO(3)$) and
the structure constants are
\[
C_{\gamma\alpha\beta}=g\epsilon_{\gamma\alpha\beta}
\]
with $\alpha$, $\beta$, and $\gamma$ running over the values 1, 2, 3.
Comparing with Eq.~\eqref{eq:17.5.33}, we see that here
\[
C_1=2.
\]
Also in this theory the matter field forms a doublet with $t_\alpha$ given by
$g/2$ times the usual Pauli matrix $\sigma_\alpha$, so
\[
C_2=1/2.
\]
Somewhat more generally, for the group $SU(N)$ with $n_f$ fermions in the
defining representation, with a conventional normalization of generators,
we have\footnote{For $SU(3)$ the $t_\alpha$ are taken as $g/2$ times the
Gell-Mann matrices $\lambda_\alpha$ used in Section 19.7, so that
$C_{\alpha\beta\gamma}=(g/2)f_{\alpha\beta\gamma}$.}
\begin{equation}
C_1=N,\qquad C_2=n_f/2.
\label{eq:17.5.35}
\end{equation}
Returning now to the general case, Eqs.~\eqref{eq:17.5.33} and
\eqref{eq:17.5.34} give
\begin{equation}
\Gamma_{A^4}^{(1\ \mathrm{loop})}
=\frac{-ig^2\mathcal I}{(2\pi)^4}\int d^4x\,F_{\gamma\mu\nu}F_\gamma^{\mu\nu}
\left[\frac{11}{12}C_1-\frac13C_2\right].
\label{eq:17.5.36}
\end{equation}
That is, the infinite constant $L_A$ in Eq.~\eqref{eq:17.4.30} is
\begin{equation}
L_A=\frac{4ig^2\mathcal I}{(2\pi)^4}
\left(\frac{11}{12}C_1-\frac13C_2\right).
\label{eq:17.5.37}
\end{equation}

It remains to say a word about the interpretation of the divergent integral
$\mathcal I$. First, before we try to integrate over the three-momentum
$\mathbf q$, we can rotate the contour of integration of $q^0$ in
Eq.~\eqref{eq:17.5.24} to the imaginary axis; as usual, the $-i\epsilon$ in
the denominator forces us to rotate counterclockwise, so that $q^0=iq^4$,
with $q^4$ running from $-\infty$ to $+\infty$. The integral is then
\begin{equation}
\mathcal I=i\int_0^\infty\frac{2\pi^2q^3\,dq}{q^4},
\label{eq:17.5.38}
\end{equation}
where $q$ is now the magnitude of the Euclidean four-vector
$(q^1,q^2,q^3,q^4)$. To go further, we evidently need some method of
regulating the integral. The simplest way of dealing with the ultraviolet
divergence is just to cut off the integral for $q$ above a scale $\Lambda$.
However, we also need a lower cut-off to deal with the infrared divergence.
This is provided by the physics of the situation. If the momenta of the four
vector particles is not zero, then a momentum flows through the internal
lines of the diagrams, providing an infrared cut-off at the scale $\mu$ of
these momenta. Similarly, if we evaluate the fourth variational derivative
of $\Gamma[A]$ with respect to $A$ not at $A=0$ but at a finite $A$, then the
propagators of the internal lines do not blow up at zero momenta, and we
have an infrared cut-off at a scale $\mu\sim gA$. Either way, $\mathcal I$
takes the form
\begin{equation}
\mathcal I=2\pi^2i\int_\mu^\Lambda\frac{dq}{q}
=2\pi^2i\ln\left(\frac{\Lambda}{\mu}\right)
\label{eq:17.5.39}
\end{equation}
and hence
\begin{equation}
L_A=-\frac{g^2}{2\pi^2}
\left(\frac{11}{12}C_1-\frac13C_2\right)
\ln\left(\frac{\Lambda}{\mu}\right)+O(g^4).
\label{eq:17.5.40}
\end{equation}
Eq.~\eqref{eq:17.4.48} then gives the renormalized coupling as
\begin{equation}
g_R=g\left[1+\frac{g^2}{4\pi^2}\ln\left(\frac{\Lambda}{\mu}\right)
\left(\frac{11}{12}C_1-\frac13C_2\right)+O(g^4)\right].
\label{eq:17.5.41}
\end{equation}
We note that while in quantum electrodynamics the radiative corrections
discussed in Section 11.2 decrease the physical coupling $g_R$ relative to
the bare coupling $g$, in non-Abelian gauge theories they \emph{increase} the
physical coupling over the bare coupling, provided that the fermion multiplet
is small enough so that $C_2<11C_1/4$. The importance of this point will be
explored in Chapter 18.

Alternatively, we can deal with the ultraviolet divergence by the methods
of dimensional regularization, discussed in Section 11.2. Here, in place of
Eq.~\eqref{eq:17.5.39}, we write
\begin{equation}
\mathcal I=i\int_0^\infty\frac{2\pi^2q^{d-1}\,dq}{(q^2+\mu^2)^2},
\label{eq:17.5.42}
\end{equation}
where $d$ is a complex dimensionality, allowed to approach 4 at the end of
the calculation, and $\mu$ is an infrared cut-off, again taken of the order of
the external momenta (or of the background fields times $g$). As long as
$d$ is complex with $\operatorname{Re}d<0$ and $\mu^2>0$, this has the finite
value
\[
\mathcal I=-i\pi^2\left(\frac d2-1\right)\mu^{d-4}
\frac{\pi}{\sin[(d/2-2)\pi]}.
\]
Analytically continuing to $d\to4$, this is
\begin{equation}
\mathcal I\longrightarrow-2i\pi^2
\left[\frac1{d-4}+\ln\mu+\cdots\right],
\label{eq:17.5.43}
\end{equation}
where $\cdots$ denotes finite $\mu$-independent terms. Here we have
\[
L_A=\frac{g^2}{2\pi^2}\left(\frac{11}{12}C_1-\frac13C_2\right)
\left(\frac1{d-4}+\ln\mu+\cdots\right)+O(g^4)
\]
and so
\begin{equation}
g_R=g\left[1-\frac{g^2}{4\pi^2}
\left(\frac{11}{12}C_1-\frac13C_2\right)
\left(\frac1{d-4}+\ln\mu+\cdots\right)+O(g^4)\right].
\label{eq:17.5.44}
\end{equation}
Note that the ultraviolet divergence here takes a different form, but the
dependence on the infrared cut-off $\mu$ is the same. Eq.~\eqref{eq:17.5.44}
will provide an important input to our discussion of asymptotic freedom in
Section 18.7.
